\section{Discussion}
\label{sec:discussion}

\subsection{Interpreting finite-size fluctuations}
\label{sec:structural-interpretation}
% Relate observed differences to joint support and community dependence.
% Answer what limiting degree matching controls and what it leaves uncontrolled:
% clustering is not isolated from all other correlations by this comparison.
% Bring together theorem-backed CM results, triangle predictions, and measured
% finite-size evidence without treating them as equally rigorous statements.
% Bring the three questions together: window scale, conditional fluctuation
% scale, and full distributional prediction. Explain where the approximation
% succeeds or fails and how that differs between models and graph conventions.
% Possible extra spread: membership fluctuations, matching variation, projection
% cascades, or omitted finite-size corrections. Treat explanations as hypotheses
% unless controlled evidence distinguishes them. Do not claim universality.
\placeholder{[Scaling evidence, approximation accuracy, and unexplained spread.]}

\subsection{Limitations and possible extensions}
\label{sec:limitations-and-extensions}
% Discuss finite sizes, grids, positive counts, allocation, calibration error,
% cutoff selection, model fit, and omitted higher-order terms. Distinguish
% evidence against a leading approximation from evidence against the exponents.
% Local weak convergence alone does not justify exchanging the
% graph-size and infinite-pruning-depth limits for the terminal core.
% Explain why the CM proof does not immediately cover the triangle ensemble.
% Possible future work: full factor-graph warning propagation, proof conditions,
% or larger communities. These are not completed experiments in this project.
% For K4 mixtures, separate automatic local 3-cores from a global transition.
\placeholder{[Finite-size and theoretical limitations, then scoped future work.]}
